%% Prospective, append-only execution addendum. It contains no scored result.
\documentclass{article}
\usepackage{amsmath}
\usepackage{booktabs}
\usepackage{hyperref}
\usepackage{seqsplit}
\usepackage{xcolor}
\usepackage{icml2024}
\makeatletter
\gdef\isaccepted{1}
\renewcommand{\Notice@String}{\textit{Internal OPHIS preregistration in ICML
two-column format. Not peer reviewed, submitted to, or accepted by ICML.}}
\makeatother
\hypersetup{
  pdfauthor={GPT Research Coordinator and independent adversarial critic},
  pdfsubject={Internal OPHIS preregistration; not an ICML submission or acceptance}
}

\newcommand{\bpb}{\mathrm{bpb}}
\newcommand{\valbpb}{\mathrm{val\_bpb}}
\newcommand{\passx}{\textcolor{blue}{pass}}
\newcommand{\killx}{\textcolor{red}{kill}}
\newcommand{\holdx}{\textcolor{orange}{hold}}
\newcommand{\hashx}[1]{{\ttfamily\seqsplit{#1}}}
\icmltitlerunning{Paper 019a: Boundary-Sidecar Attribution Addendum}

\begin{document}
\twocolumn[
\icmltitle{Paper 019a: Observe Before Optimizing\\
An Attribution Addendum for Packer-Sourced FA3 Boundaries}
\begin{icmlauthorlist}
\icmlauthor{GPT Research Coordinator}{ophis}
\icmlauthor{Systems Proposer and Independent Adversarial Critic}{ophis}
\end{icmlauthorlist}
\icmlaffiliation{ophis}{OPHIS governed autoresearch campaign}
\icmlcorrespondingauthor{GPT Research Coordinator}{baiyuzhu@mit.edu}
\icmlkeywords{automated research, causal attribution, FlashAttention, systems optimization, preregistration}
\vskip 0.3in]
\printAffiliationsAndNotice{}

\begin{abstract}
This prospective addendum strengthens only the Round-1 execution contract of
paper 019 before any governed performance diagnostic.  The candidate moves
document-boundary discovery from a per-step CUDA scan to exact metadata
already known by the packer.  A newer independent critic lowers its local
impact and transfer scores, from the paper-canonical 33/40 to 30/40, while
retaining it as a diagnostic-first dependency.  We freeze a 1,000-batch
pre-clock equivalence gate, a sequential 20-step scanner/scanner A/A phase, a
20-step scanner/sidecar A/B phase, a 10-step washout, and exactly 200 clean
timing samples.  Model, optimizer, data, boundary, loss, CPU-RNG, CUDA-RNG,
and compiler identities are bound at the relevant phase boundaries.  Two
independent H200 executions already pass adversarial event-lifetime, dynamic
shape, and real FlashAttention-3 forward/backward tests; that evidence
validates the implementation substrate but is neither a speed result nor a
RunRecord.  A single physical placement can only kill the mechanism or
require a separately preregistered role-swapped profile.  It cannot authorize
a validation-BPB endpoint.  The historical 0.927183 remains the chart origin
and a quarantined provenance observation; this addendum reports no new SOTA.
\end{abstract}

\section{Authority, Scope, and Non-Claims}

Paper 019 preregistered a connected five-round program
\cite{paper019}.  Its immutable TeX and PDF SHA-256 values are respectively
\hashx{9ba9d7bcdf7a29612bf2bc4fb6d7b8cfe69c53969187793e8f5b7960de164bc5}
and
\hashx{17a8e9d3e2af669bb4b0c105c1ec7965ad0cd9f8bc59fc7467230526db839db8}.
This append-only addendum supersedes its Round-1 execution details where the
two disagree; it does not change the other four rounds, alter their order, or
authorize an experiment.  In particular:
\[
\boxed{\begin{gathered}
\text{no scored endpoint, no validation BPB,}\\
\text{no comparator rebind, and no new SOTA.}
\end{gathered}}
\]
The chart must continue to begin at historical experiment 501, its original
timestamp, and $\valbpb=0.927183$.  That point is not a contemporaneous
control.  The active causal question remains sidecar minus scanner under the
same recovered recipe and a fresh governed seed.

This document was frozen after implementation-level systems validation but
before the governed phased diagnostic.  Preflight failures and fixes are
preserved rather than erased.  A diagnostic may execute only after typed
registration, setup reconciliation, source-hash matching, strict repository
audit, and a fresh tenant-free GPU inventory.

\section{Explicit Ratings and Selection Rationale}

Every vector uses the fixed order novelty (N), primary-source provenance (P),
causal validity (V), expected local impact (I), transfer reliability (R),
feasibility/cost (F), falsifiability (X), and connected-program coherence
(C).  All dimensions are scored from 1 to 5, where higher is better.

\begin{table}[t]
\caption{Explicit ratings frozen before the governed diagnostic.}
\label{tab:ratings}
\centering\scriptsize
\setlength{\tabcolsep}{2.2pt}
\begin{tabular}{p{0.36\linewidth}p{0.43\linewidth}r}
\toprule
assessment & N/P/V/I/R/F/X/C & sum\\
\midrule
original proposer &
4/4/5/3/4/4/5/5 & 34\\
paper-019 critic &
4/4/5/3/3/4/5/5 & 33\\
new independent critic &
4/4/5/2/2/3/5/5 & 30\\
\bottomrule
\end{tabular}
\end{table}

The lower live score is authoritative for current risk assessment.  The
candidate is not expected to improve modeling quality directly, and the
cost removed may be too small relative to GPU-placement variance.  It remains
first for a dependency reason, not because 30/40 is the largest portfolio
score: exact boundary metadata is needed to observe and attribute later
document-local mechanisms, and an accepted systems gain would increase the
charged update exposure of every downstream round.  This preserves the
campaign principle:
\begin{quote}
\emph{Before a system can improve itself, it must be able to observe and
attribute its failures.}
\end{quote}

The strongest model reserve proposed after the plateau is GPAS residual
activation control, proposer score
4/5/5/4/4/5/5/5 = 37/40 \cite{gpas2025}.  The strongest independently
criticized Engram child is distinct prime moduli per hash head/layer,
4/5/5/3/3/5/5/5 = 35/40 \cite{engram2026}.  Neither is authorized here.
They remain analysis-only successors so that a higher-scoring idea cannot
silently bypass the already frozen causal dependency.

\section{Mechanism and Plausibility}

The scanner reconstructs boundaries from the packed token tensor on every
training microbatch: it compares BOS locations, performs data-dependent
compaction, and materializes cumulative sequence lengths on CUDA.  The
packer already knows every row start and interior document start while it
writes the batch.  The treatment records those flattened offsets into a
three-slot pinned-CPU \texttt{int32} ring, asynchronously copies only the
logical prefix to CUDA, and protects reuse with producer and consumer events.
The same offset tensor supplies FlashAttention-3's
\texttt{cu\_seqlens\_q} and \texttt{cu\_seqlens\_k}.

The causal path is therefore
\[
\begin{aligned}
\text{packer-known starts}
&\rightarrow \text{exact pinned sidecar}\\
&\rightarrow \text{remove CUDA boundary scan}\\
&\rightarrow \text{lower charged step time}\\
&\rightarrow \text{more updates in 300 seconds}\\
&\rightarrow \text{possibly lower endpoint BPB}.
\end{aligned}
\]
The mechanism is plausible only if every semantic and state mediator is
identical and the timing gain clears a robust break-even.  The FA3 source
supports the external cumulative-boundary interface \cite{fa3}; it does not
support the local speed or BPB claim.  That claim must come entirely from
the governed local contrast.

\section{Frozen Phased Diagnostic}

\paragraph{Symmetric pre-clock assay.}
Both registered arms independently instantiate scanner and sidecar loaders
and compare exactly 1,000 batches before the diagnostic clock.  Inputs,
targets, epochs, flattened cumulative boundaries, boundary counts, maximum
segment length, and SHA-256 digests must match.  Required invariants are zero
first offset, terminal offset $BT$, strict increase, every row start, and no
gap greater than $T$.  An interior-BOS mismatch is a hard kill.

\paragraph{Steps 0--19: A/A attribution prefix.}
Both physical arms consume scanner-derived masks.  The treatment advances and
safely discards its sidecar lease so producer behavior remains exercised; this
conservative asymmetry is disclosed and is not used as a speed sample.
Immediately after update 20, exact data, boundary, micro-loss, model,
optimizer, CPU-RNG, and CUDA-RNG hashes must agree across arms.

\paragraph{Steps 20--39: A/B mechanism phase.}
At step 20, the treatment alone switches its consumed mask metadata to the
sidecar; the control remains on the scanner.  After update 40, the same
phase-local data, boundary, loss, state, and RNG hashes must agree.  Any state
divergence kills the mechanism, even if timing appears favorable.

\paragraph{Steps 40--49: washout.}
These ten steps are executed but excluded from the performance estimator.
They prevent the activation boundary, Python allocation transients, and
phase-local instrumentation from entering the clean window.

\paragraph{Steps 50--249: clean profile.}
Exactly 200 positive finite step-time observations are retained.  They form
20 consecutive blocks of 10, and the same complete blocks determine both the
point estimate and its one-sided 95\% lower confidence bound.  The bootstrap
uses the frozen seed 19019063 and 10,000 resamples.  No in-clock state hashing,
boundary hashing, or graph probing is admitted to this window unless charged
symmetrically.

\paragraph{Compiler identity.}
The first model call may create exactly one PyTorch Dynamo graph at site
\texttt{0:0}.  Every later model call executes under
\texttt{fail\_on\_recompile}.  Callback-derived compile IDs are hashed; zero
recompiles and zero CUDA-graph recordings are required.  Unsupported model or
optimizer state types fail closed rather than falling back to
representation-dependent hashing.

\section{Two-Level Performance Decision}

For the first physical placement, the treatment/control token-rate point
ratio must be at least 1.054 and its one-sided 95\% lower bound at least
1.041.  The point threshold rounds up the robust scientific break-even from
paper 019,
\[
r_{\rm robust}
=\exp(0.002396/0.046)=1.05347.
\]
The separate 1.041 lower-bound screen approximates the central, rather than
robust, token-law break-even and is a fixed diagnostic decision bound rather
than population-level experimental confidence.
Failure of either preregistered diagnostic threshold kills Round 1 before
validation scoring.

Passing both thresholds does \emph{not} authorize a 300-second endpoint.
One GPU may be systematically faster than the other, and temporal
bootstrapping cannot identify placement bias.  A pass yields only
\texttt{REQUIRE\_COUNTERBALANCED\_PROFILE}.  A separately frozen role-swapped
AB/BA profile must reuse the same source, challenge, UUID bindings, phase
contract, and complete-block estimator.  Only a counterbalanced pass may
promote the ordinary two-arm endpoint proposal through the normal typed gate.
Thus a single placement can kill or request stronger evidence, but cannot
create a favorable BPB opportunity.

Each pair artifact must bind exactly two distinct immutable RunRecords, their
fingerprints, roles, run UUIDs, source snapshot, challenge fingerprint, and
resolved configurations.  The sole resolved arm difference is
\texttt{PACKER\_DOC\_BOUNDARIES}.  A child failure still writes an
\texttt{INVALID\_DIAGNOSTIC} pair artifact.  Recovery may finalize comparison
from existing records but may never relaunch the same immutable diagnostic.

\section{H200 Systems Evidence Before Execution}

Two independent non-experimental validators ran on H200s: one with ordinary
asynchronous CUDA execution and one with
\texttt{CUDA\_LAUNCH\_BLOCKING=1}.  Both produced the same terminal digest and
passed delayed H2D, delayed consumer work across ring wraps, outstanding
fourth-lease rejection, wrong-stream rejection, discard contracts, poisoned
tail lengths, exact capacity bounds, a single dynamic compile identity,
forced static-shape recompile rejection, and real FA3 varlen
forward/backward bitwise equivalence over nine ring-wrap iterations.
The FA3 check used a synthetic 32-token, two-head, head-dimension-64 workload
and a small compiled helper.  It did not execute the full training graph.

\begin{table}[t]
\caption{Preflight evidence.  ``Pass'' is systems scope only.}
\label{tab:h200}
\centering\scriptsize
\setlength{\tabcolsep}{2pt}
\begin{tabular}{p{0.49\linewidth}p{0.26\linewidth}}
\toprule
check & result\\
\midrule
runtime & PyTorch 2.9.1+cu128; CUDA 12.8\\
device & NVIDIA H200, SM90\\
event and ring lifetime & \passx\\
Dynamo graphs / recompiles & 1 / 0\\
forced recompile & rejected\\
real FA3 forward/backward & bitwise exact\\
async vs blocking log SHA & identical\\
\bottomrule
\end{tabular}
\end{table}

Two predecessor attempts are part of the record.  The first failed before GPU
work because the standalone validator lacked the repository root on
\texttt{sys.path}; the second found that applying
\texttt{fail\_on\_recompile} to the initial call rejects the required first
compile.  The frozen implementation now permits the first compile and applies
the stance to all later calls.  These failures improve attribution; they are
not discarded failed trials.

The validated FA3 result digest is
\hashx{e5a96b00475fbcf1425eb72315ef0ef02e6138fc82cb84983529ef452f750127}.
The two log files share SHA-256
\hashx{37bae79b44198818fce02ef54a6a26482f119505511a547fa9fd6a17759ac906}.
This evidence is neither a RunRecord nor a throughput measurement.

\section{Retrospective Round-2 Update}

Paper 019's canonical-key Round 2 was independently evaluated by its frozen
offline mediator before any GPU scoring.  The 8,192-token vocabulary reduced
to 6,808 canonical keys (16.8945\%), while the conservative observed affected
shares were 17.7887\% for bigrams and 3.7155\% for trigrams.  The governing
minimum, 3.7155\%, was below the preregistered 10\% gate.  The result is
\[
\boxed{\texttt{KILL\_BEFORE\_GPU\_SCORING}.}
\]
The result and output-manifest SHA-256 values are
\hashx{9d0e3a122a7e880471bc793368576a89afa5cfbe442f1149c1bbc651eed4d761}
and
\hashx{d15dc6b3be202b463befd270abf1af04b8b24c3ab3a6acce831a7a5a58189f63}.
Pooled or less conservative shares cannot replace the frozen statistic.
This result was known only after paper 019 and may update the portfolio, but
it may not tune Round-1 thresholds.  A future paper may replace the dead round
with the independently ranked prime-modulus $\rightarrow$ constant-parameter
$K=4$ $\rightarrow$ contextual-gate Engram chain.

\section{Reproducibility Binding}

The exact systems-validation manifest is stored under
\path{research/experiments/artifacts/paper019_round1_fa3_boundary_sidecar}.
At this addendum freeze, the principal source SHA-256 values are:
\begin{center}
\scriptsize
\begin{tabular}{p{0.95\linewidth}}
lib.py\\
\hashx{13de28bcb398866be65cfca489413818d75dc2fcd2461e367be06198f879123d}\\[2pt]
train.py\\
\hashx{e2d4ac9e91c6ef716a8fdab35ab72da3e0e4bb23cbbecee04cc528b66d172b8b}\\[2pt]
tools/run\_gated.py\\
\hashx{47d08ec38f6dff53ed50cdc5e4ddaf88342da0f2e9c5993219baf4121711868a}\\[2pt]
tools/run\_stage.py\\
\hashx{c7d3ab45af2983598936859ea9f14c34ebb6da3ce9b1d0cc9d00ca6c764b16c0}\\[2pt]
tools/register\_paper019\_fa3\_boundary\_sidecar\_chain.py\\
\hashx{d8f6b4e71a65c971ceffdcb63630b99daa4d4cd57c3546cc52410558ee8e9ce0}\\[2pt]
tools/validate\_packer\_sidecar\_h200.py\\
\hashx{e5d58f8d0a3db1a87b1e82ee1b4063da92e49eac015a7c6b6e0f720379a1c9e4}\\[2pt]
H200 validation manifest\\
\hashx{f0a226c34198186bfce9a047e328d2b4e3fb290a3993dc1cb24829792b37c213}
\end{tabular}
\end{center}
Any later edit invalidates these bindings and requires a new reconciliation
and explicit amendment before execution.  Remote launch must independently
match the local source hashes and the pinned FA3 repository.

\section{Conclusion}

This addendum converts an attractive exact-semantics optimization into a
failure-attributable sequence.  It separates semantic equivalence, state
equivalence, compiler identity, systems lifetime, clean timing, placement
bias, and endpoint BPB.  The candidate can advance only by clearing each
link in order; a favorable downstream number cannot repair an upstream
failure.  That discipline is the mechanism by which exploration remains
novel without becoming ungrounded.

{\footnotesize
\begin{thebibliography}{9}

\bibitem[GPT Research Coordinator(2026)]{paper019}
GPT Research Coordinator.
Observe, Attribute, Then Improve: A Critic-Ranked Five-Mechanism Program for
an Autoresearch Plateau. OPHIS paper 019, 2026.

\bibitem[Dao et al.(2024)]{fa3}
T. Dao et al.
FlashAttention-3: Fast and Accurate Attention with Asynchrony and
Low-precision. NeurIPS, 2024.

\bibitem[Chen et al.(2025)]{gpas2025}
T. Chen et al.
GPAS: Accelerating Convergence of LLM Pretraining via
Gradient-Preserving Activation Scaling.
NeurIPS, 2025.

\bibitem[Cheng et al.(2026)]{engram2026}
X. Cheng et al.
Conditional Memory via Scalable Lookup: A New Axis of Sparsity for Large
Language Models. ACL, 2026.

\end{thebibliography}
}

\end{document}
